%% Chapter 6: Conclusion and Future Work: chapters/chapter06.tex

% ==============================================================================
% CONTENT BLUEPRINT — Chapter 6: Conclusion and Future Work (M.S. Thesis)
% ==============================================================================
% Reading audience: anyone reading the thesis to form a verdict on it.
% Goal: in 6–10 pages, deliver the verdict on every objective (§1.4) and
% roadmap future work.
%
% ----------------------------------------------------------------------------
% §6.1 Summary of Findings  (2–3 pages)
%   • Open with a three-sentence verdict: what was studied, how, and what
%     was found (with one headline number).
%   • Bullet list of headline findings, one per §1.4 objective, each
%     referencing the supporting Chapter 4/5 subsection.
%   • Explicit verdict on each objective: met / partially met / not met,
%     with reason.
%
% §6.2 Major Contributions  (½–1 page)
%   • Enumerate:
%       – Code / scripts deposited (GitHub repo + Zenodo DOI),
%       – Datasets deposited (Zenodo DOI),
%       – Calibration / parameter tables,
%       – Methodological improvements (workflow, SOPs, automation),
%       – Any open-source tools contributed back upstream.
%   • This is what can be reused by the next scholar.
%
% §6.3 Limitations of the Work  (½–1 page)
%   • Restate quantitative limitations from §4.7 and §5.8.
%   • Identify irreducible approximations (zero-temperature, particular
%     functional, neglect of phonons, etc.).
%   • Be concise — one sentence per limitation is typical.
%
% §6.4 Future Research Directions  (1–2 pages)
%   • 3–5 specific, actionable directions:
%       – What would be done (one sentence)
%       – What new physics it unlocks (one sentence)
%       – What is required (compute, software, time, collaborations).
%   • Each must be specific enough for another student to pick up.
%
% §6.5 Closing Reflection  (¼ page, optional)
%   • One paragraph reflecting on the broader significance of the work or
%     the candidate's intellectual trajectory. Avoid family / feelings.
%
% Writing-style rules:
%   • Use declarative past tense ("we found", "the calculations show").
%   • Each numerical claim must point to §4 or §5 (or §1 if repeating).
%   • No new figures. If you find you need one, go back to Chapter 4 or 5.
%   • If a future-work item is essentially "do this thesis again with X",
%     it is weak — re-cast as a specific physics question.
% ==============================================================================

\chapter{Conclusion and Future Work}
\label{ch:conclusion}

\section{Summary of Findings}
In this thesis, we performed a systematic first-principles study on the structural, electronic, and optical properties of transition metal dichalcogenide monolayers. The key findings are:
\begin{enumerate}
    \item Structural relaxation confirmed stable geometries, with lattice constants agreeing with experimental values within 1.5\%.
    \item Biaxial strain was shown to tune the electronic bandgap linearly up to $\pm 3\%$ strain, with a transition from direct to indirect bandgap under higher tensile strain.
    \item The calculated optical absorption coefficients show strong absorption exceeding $10^5$ cm$^{-1}$ in the visible energy range, which is enhanced under strain.
\end{enumerate}

\section{Major Contributions}
This research contributes to the fundamental understanding of strain-engineered two-dimensional semiconductors. The detailed calculations of carrier effective masses and optical constants provide useful quantitative parameters for designing strain-tunable optoelectronic devices.

\section{Limitations of the Work}
The primary limitations of this study are:
\begin{itemize}
    \item All calculations were performed at $0$ K temperature, ignoring thermal expansion and electron-phonon scattering effects.
    \item Standard DFT functionals were used, which do not completely account for excitonic effects on the optical absorption spectra.
\end{itemize}

\section{Future Research Directions}
For future investigations, we propose:
\begin{itemize}
    \item Incorporating temperature effects using ab-initio molecular dynamics (AIMD) simulations.
    \item Calculating excitonic properties using the Bethe-Salpeter Equation (BSE) on top of $GW$ calculations to obtain highly accurate optical absorption profiles.
    \item Investigating the electronic transport properties under strain using BoltzTraP code.
\end{itemize}
